\documentclass[12pt]{article}

\usepackage[a4paper,margin=1in]{geometry}
\usepackage[T1]{fontenc}
\usepackage{longtable}
\usepackage{booktabs}
\usepackage{array}
\usepackage{enumitem}
\usepackage{hyperref}

\setlength{\parskip}{0.6em}
\setlength{\parindent}{0pt}
\setlist[itemize]{leftmargin=1.6em}

\title{Integration Testing Report for StreetRace Manager}
\author{Lohith Kola}
\date{\today}

\begin{document}

\maketitle

\section{Introduction}
For Part 2, I built a small command-line system called StreetRace Manager and focused on how the modules interact with each other. Instead of testing each module in isolation, I designed the system so that registration, crew roles, inventory, races, results, missions, garage work, and rankings all share state through one integration-friendly facade. The main goal of this part was to verify data flow, business rules, and state updates across module boundaries.

\section{Codebase Overview}
The system is implemented in Python under the \texttt{integration/code} folder. I kept the design in-memory so the integration behavior is easy to trace and test.

\subsection{Modules Implemented}
\begin{itemize}
  \item \textbf{Registration module}: registers new crew members and stores their names.
  \item \textbf{Crew Management module}: assigns roles such as driver, mechanic, and strategist, and stores skill levels.
  \item \textbf{Inventory module}: tracks cars, spare parts, tools, and cash balance.
  \item \textbf{Race Management module}: creates races and validates driver/car entry.
  \item \textbf{Results module}: records outcomes, updates cash, updates rankings, and marks vehicles damaged when needed.
  \item \textbf{Mission Planning module}: creates missions, checks required roles, and validates mission startup.
  \item \textbf{Garage module}: handles vehicle inspection and repair flow.
  \item \textbf{Rankings module}: maintains the leaderboard, points, wins, and reputation.
\end{itemize}

\subsection{Why I Added Two Extra Modules}
I added the Garage module and the Rankings module because both of them make the integration part richer and more realistic.

\begin{itemize}
  \item \textbf{Garage}: This made it possible to connect race damage, mechanic availability, spare parts, tools, and repair completion into one flow.
  \item \textbf{Rankings}: This gave the Results module a meaningful downstream effect beyond just storing a race result. Recording a race now updates the leaderboard and crew reputation.
\end{itemize}

\subsection{System Facade}
I used \texttt{system.py} as the integration facade. This file wires all modules together and gives the CLI and tests one consistent entry point. That made the integration tests much cleaner because they could focus on business flows instead of manually connecting every service object each time.

\section{Control Flow Graph}
I generated a detailed textual call graph in \texttt{integration/diagrams/call\_graph\_instructions.md}. That file lists the main functions in each module, the cross-module calls, and an ASCII call graph that is easy to copy into a hand-drawn version.

I will draw the final call graph by hand and attach the image separately in the final submission.

\section{Integration and Business Rules}
The implementation enforces the following integration rules:

\begin{itemize}
  \item A crew member must be registered before any role can be assigned.
  \item Only crew members with the driver role may be entered into a race.
  \item A race entry also requires an available and undamaged car from inventory.
  \item Recording a race result updates the inventory cash balance.
  \item Recording a race result also updates the rankings module.
  \item Missions cannot start unless the assigned crew cover every required role.
  \item If a damaged car is involved in a repair flow, a mechanic must be available before the mission can proceed.
  \item Repair work checks tool availability and spare-part availability before restoring the car.
\end{itemize}

\section{Integration Test Design}
\subsection{Approach}
I wrote the integration tests to follow the real business flows in the system instead of only checking single methods. Each test goes through more than one module and verifies that the state changes stay consistent across them.

\subsection{Tests Implemented}
\begin{longtable}{p{0.21\textwidth}p{0.18\textwidth}p{0.19\textwidth}p{0.19\textwidth}p{0.19\textwidth}}
\toprule
\textbf{Test} & \textbf{Modules involved} & \textbf{Expected result} & \textbf{Actual result} & \textbf{Why the test matters} \\
\midrule
\endhead
\texttt{test\_registering\_driver\_and\_entering\_race\_succeeds} & Registration, Crew Management, Inventory, Race Management & Registered driver enters race with a valid car and race state becomes ready & Passed; driver, car, and race state updated correctly & This is the basic integration path that proves the core system can actually create a legal race entry. \\
\texttt{test\_entering\_race\_with\_unregistered\_member\_fails} & Registration, Race Management & Entry should be rejected with an exception & Passed; invalid entry raised the correct failure & This checks the rule that registration must happen before race participation. \\
\texttt{test\_entering\_race\_with\_registered\_non\_driver\_fails} & Registration, Crew Management, Race Management & Registered non-driver should still be rejected & Passed; driver-role validation worked & This is needed because registration alone is not enough for race eligibility. \\
\texttt{test\_recording\_race\_result\_updates\_cash\_and\_rankings} & Race Management, Results, Inventory, Rankings & Prize money should raise cash and rankings should update & Passed; cash, wins, points, leaderboard, and race status all updated & This checks whether race completion triggers the required downstream state changes. \\
\texttt{test\_recording\_damaged\_race\_result\_keeps\_car\_unavailable\_for\_next\_race} & Results, Inventory & Damaged car should remain unavailable after result recording & Passed; damaged car stayed unavailable & This is important because a damaged car should not accidentally be reused immediately. \\
\texttt{test\_mission\_cannot\_start\_when\_required\_roles\_are\_unavailable} & Registration, Crew Management, Mission Planning & Mission start should fail when one required role is missing & Passed; mission raised the expected error & This verifies the main mission availability rule. \\
\texttt{test\_mission\_starts\_when\_required\_roles\_are\_present} & Registration, Crew Management, Mission Planning & Mission should start once the required roles are covered & Passed; mission moved to started state & This confirms the positive branch for role validation. \\
\texttt{test\_damaged\_car\_requires\_mechanic\_before\_repair\_mission\_starts} & Results, Inventory, Mission Planning & Repair mission should fail without mechanic coverage & Passed; mechanic requirement was enforced & This directly tests the damaged-car business rule from the assignment. \\
\texttt{test\_repair\_flow\_consumes\_parts\_and\_restores\_car} & Results, Mission Planning, Garage, Inventory & Parts should be consumed and the car should return to ready state & Passed; part count dropped and the car was repaired & This verifies the repair path across several modules, not just the garage function alone. \\
\texttt{test\_repair\_fails\_when\_non\_mechanic\_attempts\_garage\_work} & Crew Management, Garage & Garage repair should reject non-mechanics & Passed; invalid repair raised the expected exception & This prevents role bypass in the repair path. \\
\texttt{test\_end\_to\_end\_flow\_across\_registration\_race\_results\_and\_missions} & Registration, Crew Management, Inventory, Race Management, Results, Mission Planning, Garage, Rankings & Full scenario should work from setup to race, damage, repair, and second race entry & Passed; the whole scenario completed correctly & This gives confidence that the system works as one integrated application instead of just separate modules. \\
\bottomrule
\end{longtable}

\subsection{Errors or Logical Issues Found}
The final integration suite passed on the first full run after I finished the module implementation, so the test phase did not expose any extra post-implementation defects that required a separate debug cycle. In other words, once the full integrated design was in place, all 11 integration scenarios behaved as expected.

\section{Test Execution Summary}
I ran the integration suite with:

\begin{itemize}
  \item \texttt{python3 -m pytest integration/tests -v}
\end{itemize}

Final result:

\begin{itemize}
  \item \textbf{11 tests passed}
  \item \textbf{0 tests failed}
\end{itemize}

The tests covered the main required interactions:

\begin{itemize}
  \item registration to role assignment,
  \item role validation during race entry,
  \item result recording to inventory/rankings update,
  \item role validation during mission startup,
  \item damaged-car repair flow with mechanic requirement,
  \item one full end-to-end scenario across most modules.
\end{itemize}

\section{Errors Found and Fixes Applied}
No additional errors were discovered during the final integration test run. Because the test suite passed cleanly, no extra bug-fix commit was needed after the integration tests were introduced.

\section{How to Run the Code and Tests}
From the repository root, I can run the demo CLI with:

\begin{itemize}
  \item \texttt{python3 -m integration.code.cli}
\end{itemize}

I can run the integration tests with:

\begin{itemize}
  \item \texttt{python3 -m pytest integration/tests -v}
\end{itemize}

\section{Hand-Drawn Call Graph Note}
The assignment requires a hand-drawn call graph image. I therefore created a detailed textual guide in \texttt{integration/diagrams/call\_graph\_instructions.md} and a reminder file in \texttt{integration/diagrams/hand\_drawn\_call\_graph\_placeholder.txt}. I should use those instructions to draw the graph by hand, take a clear photo, and include that image in the final submission.

\section{Final Reflection}
This part was a good example of why integration testing matters. Most of the interesting behavior in StreetRace Manager is not inside a single method, but in the way one module hands state to the next one. A race result changing cash and rankings, a damaged car blocking future race use, and a repair flow needing a mechanic, tools, and parts are all cross-module behaviors. Those were exactly the interactions I wanted to verify.

Keeping the system in memory also helped a lot because it made the data flow easy to inspect during testing. The final structure is simple, readable, and ready to extend if more race or mission features are added later.

\end{document}
